\uuid{SbqG}
\exo7id{6153}
\titre{exo7 6153}
\auteur{queffelec}
\organisation{exo7}
\datecreate{2011-10-16}
\isIndication{false}
\isCorrection{false}
\chapitre{Connexité}
\sousChapitre{Connexité}
\cpge{mpsi-mp2i.raisonnement-vocabulaire-ensembliste.c, pcsi.raisonnement-vocabulaire-ensembliste.d, ptsi.raisonnement-vocabulaire-ensembliste.d, mpsi-mp2i.limites-continuite.c, pcsi.limites-continuite.c, ptsi.limites-continuite.c, mpsi-mp2i.derivabilite.a, pcsi.derivabilite.a, ptsi.derivabilite.a, bcpst1.vocabulaire-applications, bcpst1.limites-continuite.c, bcpst1.derivation.a}
\module{Topologie}
\niveau{L3}
\difficulte{}

\contenu{
\texte{
Soit $I=[a,b]$ et $f : I\to\R$ dérivable~; soit $A=\big\{(x,y)\in
I\times I; y>x\big\}$  et $g : A\to \R$ définie par
$g(x,y)={f(y)-f(x)\over y-x}$.
}
\begin{enumerate}
    \item \question{Montrer que $g(A)\subset f'(I)\subset\overline {g(A)}$.}
    \item \question{Montrer que $f'$ a la propriété de la valeur intermédiaire~: si elle
prend les valeurs
$\alpha $ et $\beta $, elle prend toute valeur $\gamma \in[\alpha ,\beta ]$.}
\end{enumerate}
}
